\documentclass[12pt]{article}
\usepackage[utf8]{inputenc}
\usepackage[T1]{fontenc}
\usepackage[a4paper,margin=3cm]{geometry}
\usepackage{minted}
\usepackage{xcolor}
\usepackage{csquotes}

\setlength{\parindent}{0em}
\setlength{\parskip}{\medskipamount}

\title{Container in Python: Sammeln von Daten}
\author{}
\date{}

\begin{document}
	
	\maketitle
	
	\section*{Einleitung}
	
	Neben Listen bietet Python weitere grundlegende Datenstrukturen, die spezifische Eigenschaften aufweisen und für bestimmte Aufgaben besser geeignet sind. In diesem Kapitel behandeln wir zunächst \texttt{dict} und \texttt{set}, zwei sogenannte \emph{unstrukturierte} Datenstrukturen. Anschließend erläutern wir das Konzept von \emph{Veränderbarkeit} (\emph{Mutability}) in Python und erklären, warum die Datenstruktur \texttt{tuple} trotz der Existenz von Listen sinnvoll ist.
	
	\section*{1. Dictionaries – Schlüssel-Wert-Zuordnungen}
	
	Ein \texttt{dict} (Wörterbuch) speichert eine Zuordnung zwischen eindeutigen Schlüsseln und Werten. Die Elemente eines Dictionaries sind somit nicht einfach geordnete Daten, sondern benannte bzw. identifizierte Werte.
	
	Ein Dictionary wird mit geschweiften Klammern erstellt, wobei die Schlüssel-Wert-Paare durch Doppelpunkte getrennt sind:
	
\begin{minted}[bgcolor=gray!10]{python}
noten = {"Anna": 1.3, "Bernd": 2.0}
\end{minted}

\begin{minted}[bgcolor=gray!10]{python}
noten = {}
#oder
noten = dict()

noten["Anna"] = 1.3
noten["Bernd"] = 2.0
\end{minted}
	
	Zugriff erfolgt über die Schlüssel:
	
\begin{minted}[bgcolor=gray!10]{python}
noten["Anna"]    # ergibt 1.3
\end{minted}
	
	Einträge lassen sich entfernen:
	
\begin{minted}[bgcolor=gray!10]{python}
del noten["Bernd"]
\end{minted}
	
	Wichtige Methoden:
	
\begin{minted}[bgcolor=gray!10]{python}
noten.keys()     # alle Schlüssel
noten.values()   # alle Werte
noten.items()    # Paare als Tupel
\end{minted}
	
	Iteration über Schlüssel und Werte:
	
\begin{minted}[bgcolor=gray!10]{python}
for name, note in noten.items():
  print(name, "hat", note)
\end{minted}

	
	\section*{2. Sets – Mengen ohne Duplikate}
	
	Ein \texttt{set} ist eine ungeordnete Sammlung von einzigartigen (d.h. nicht mehrfach vorkommenden) Elementen. Die Datenstruktur ähnelt mathematischen Mengen.
	
	Erstellung mit geschweiften Klammern oder \texttt{set()}:
	
\begin{minted}[bgcolor=gray!10]{python}
zahlen = {1, 2, 3, 2}
print(zahlen)  # {1, 2, 3}
\end{minted}

	Achtung! Mit \verb|zahlen = {}| wird ein \verb|dict| erzeugt. Wenn wir eine Menge erzeugen wollen, müssen wir die geschwungenen Klammern mit Werten füllen oder den Konstruktor \verb|zahlen = set()| verwenden.

\begin{minted}[bgcolor=gray!10]{python}
zahlen = {1, 2, 3, 2}
#oder
zahlen = set()
\end{minted}
	
\begin{minted}[bgcolor=gray!10]{python}
aus_liste = set([1, 1, 2])
print(aus_liste)  # {1, 2}
\end{minted}
	
	Elemente können hinzugefügt und entfernt werden:
	
\begin{minted}[bgcolor=gray!10]{python}
zahlen.add(4)
zahlen.remove(2)
\end{minted}
	
	Typische Mengenoperationen:
	
\begin{minted}[bgcolor=gray!10]{python}
a = {1, 2, 3}
b = {3, 4, 5}

a | b   # Vereinigung: {1, 2, 3, 4, 5}
a & b   # Durchschnitt: {3}
a - b   # Differenz: {1, 2}
\end{minted}
	
	\section*{3. Veränderbarkeit: mutable vs. immutable}
	
	Python unterscheidet zwischen veränderbaren (\emph{mutable}) und unveränderlichen (\emph{immutable}) Objekten:
	
	\begin{itemize}
		\item \textbf{Mutable:} Der Inhalt des Objekts kann nachträglich verändert werden (z.\,B.\ \texttt{list}, \texttt{dict}, \texttt{set}).
		\item \textbf{Immutable:} Der Inhalt ist fest und kann nicht verändert werden (z.\,B.\ \texttt{int}, \texttt{str}, \texttt{tuple}).
	\end{itemize}
	
	Warum ist das wichtig? Nur \emph{unveränderliche Objekte} können als Schlüssel in einem Dictionary oder als Elemente eines Sets verwendet werden. Denn nur dann kann sichergestellt werden, dass sich ihr Wert (und damit ihr Hash-Wert) nicht nachträglich ändert.
	
	\section*{4. Tuple – Geordnete, unveränderliche Daten}
	
	Ein \texttt{tuple} ist eine geordnete, aber \emph{unveränderliche} Datenstruktur. Sie ähnelt der Liste, unterscheidet sich aber in zwei zentralen Punkten:
	
	\begin{itemize}
		\item Die Elemente eines Tupels können nach der Erstellung nicht verändert werden.
		\item Tupel sind \emph{hashable} – sie können als Schlüssel in Dictionaries oder als Elemente in Sets verwendet werden.
	\end{itemize}
	
	Erstellung:
	
\begin{minted}[bgcolor=gray!10]{python}
punkt = (3, 4)
paar = "Hallo", 42       # Klammern optional
eins = (7,)              # Einzelwert-Tuple
\end{minted}
	
	Zugriff per Index:
	
\begin{minted}[bgcolor=gray!10]{python}
punkt[0]    # ergibt 3
\end{minted}
	
	Veränderung ist nicht erlaubt:
	
\begin{minted}[bgcolor=gray!10]{python}
punkt[0] = 5    # führt zu einem Fehler
\end{minted}

	Verwendung als Dictionary-Schlüssel:
	
\begin{minted}[bgcolor=gray!10]{python}
messung = {}
messung[(2024, 5)] = 19.3
\end{minted}
	
	Tupel sind ideal für feste, strukturierte Datenpakete und werden häufig verwendet, um mehrere Rückgabewerte aus Funktionen zu liefern:
	
\begin{minted}[bgcolor=gray!10]{python}
def teile(x, y):
  return x // y, x % y

ganz, rest = teile(17, 5)
\end{minted}

	\section*{5. Zusammenfassung}
	
	\begin{tabular}{l|c|c|c}
		Datenstruktur & Geordnet & Veränderbar & Duplikate erlaubt \\
		\hline
		\texttt{list} & ja       & ja          & ja \\
		\texttt{tuple} & ja      & nein        & ja \\
		\texttt{set} & nein      & ja          & nein \\
		\texttt{dict} & nein     & ja          & Schlüssel eindeutig \\
	\end{tabular}
	
	Jede dieser Datenstrukturen hat ihre eigene Stärke. In der Praxis ist die Wahl der richtigen Struktur entscheidend für effizienten und verständlichen Code.
	
	
	
\end{document}
